\documentclass[11pt]{article}
\usepackage[a4paper,margin=1in]{geometry}
\usepackage[T1]{fontenc}
\usepackage[utf8]{inputenc}
\usepackage{lmodern}
\usepackage{microtype}
\usepackage{amsmath,amssymb,amsthm,mathtools}
\usepackage{xcolor}
\usepackage{hyperref}
\usepackage{tcolorbox}
\usepackage{enumitem}
\usepackage{array}
\usepackage{makecell}
\usepackage{multirow}
\usepackage{titlesec}
\usepackage{fancyhdr}

\hypersetup{colorlinks=true,linkcolor=blue!60!black,urlcolor=blue!60!black,citecolor=blue!60!black}
\setlength{\parindent}{0pt}
\setlength{\parskip}{0.6em}
\setlength{\headheight}{14pt}
\pagestyle{fancy}
\fancyhf{}
\lhead{Zero Knowledge: An Introduction}
\rhead{\thepage}

% macro.tex -- all macros for zero-knowledge notes

%% ---- tcolorbox libraries ----
\tcbuselibrary{skins}

%% ================================================================
%%  Standard notation
%% ================================================================
\newcommand{\G}{\mathbb{G}}
\newcommand{\Z}{\mathbb{Z}}
\newcommand{\F}{\mathbb{F}}
\newcommand{\Adv}{\mathcal{A}}
\newcommand{\set}[1]{\left\{#1\right\}}

% NP relation and language
\newcommand{\relation}{\mathcal{R}}
\newcommand{\lang}{\mathcal{L}}

% NIZK game macros
\newcommand{\statement}{x}
\newcommand{\zkproof}{\pi}
\newcommand{\crs}{\mathsf{crs}}
\newcommand{\td}{\mathsf{td}}
\newcommand{\simsetup}{\mathsf{SimSetup}}
\newcommand{\simprove}{\mathsf{SimProve}}
\newcommand{\verify}{\mathsf{Verify}}

% ZK-specific shorthands
\newcommand{\Zq}{\Z_q}            % Z/qZ
\newcommand{\deq}{\mathrel{:=}}   % definitional assignment

%% ================================================================
%%  Theorem-like environments
%% ================================================================
\newtheorem{definition}{Definition}
\newtheorem{remark}{Remark}
\newtheorem{example}{Example}

%% ================================================================
%%  Typography: section headings and default tcolorbox style
%% ================================================================
\titleformat{\section}{\Large\bfseries\color{blue!45!black}}{\thesection}{0.6em}{}
\titleformat{\subsection}{\large\bfseries\color{blue!35!black}}{\thesubsection}{0.6em}{}
\titleformat{\subsubsection}{\normalsize\bfseries\color{blue!30!black}}{\thesubsubsection}{0.6em}{}

\tcbset{colback=blue!2!white,colframe=blue!50!black,boxrule=0.6pt,arc=1.5mm}

%% ================================================================
%%  Colors for protocol diagrams
%% ================================================================
\definecolor{zkborder}{RGB}{35,85,175}
\definecolor{zkback}{RGB}{247,250,255}
\definecolor{zkpropborder}{RGB}{55,120,195}
\definecolor{zkpropback}{RGB}{240,247,255}

%% ================================================================
%%  Property checklist
%% ================================================================

%  Colored Yes / No indicators
\newcommand{\propyes}{\textcolor{green!65!black}{\textbf{Yes}}}
\newcommand{\propno}{\textcolor{red!65!black}{\textbf{No}}}

%  \propbox{completeness}{honest-verifier ZK}{special soundness}
\newcommand{\propbox}[3]{%
  \begin{tcolorbox}[
    enhanced,
    colback=zkpropback, colframe=zkpropborder,
    boxrule=0.6pt, arc=2.5mm,
    left=5pt, right=5pt, top=4pt, bottom=4pt,
    fontupper=\small
  ]
    \begin{itemize}[noitemsep, leftmargin=1.3em]
      \item Completeness?\quad #1
      \item Honest-verifier ZK?\quad #2
      \item Special soundness?\quad #3
    \end{itemize}
  \end{tcolorbox}%
}

%% ================================================================
%%  Protocol diagram components
%% ================================================================

%  Message arrows
\newcommand{\send}[1]{\ensuremath{\xrightarrow{\;\;#1\;\;}}}   % Prover -> Verifier
\newcommand{\recv}[1]{\ensuremath{\xleftarrow{\;\;#1\;\;}}}    % Verifier -> Prover

%  Outer frame for the protocol table.
%  Usage: wrap \begin{sigmatab}...\end{sigmatab} inside this.
\newtcolorbox{protocolbox}{
  enhanced,
  colback=zkback, colframe=zkborder,
  arc=2.5pt, boxrule=1pt,
  left=5pt, right=5pt, top=5pt, bottom=5pt
}

%  Protocol table (tabular only — place inside protocolbox).
%  Arguments: #1 = prover label, #2 = verifier label.
%  Include $...$ yourself if the label contains math.
%
%  Example:
%    \begin{protocolbox}
%    \begin{sigmatab}{$P(h;\,w)$}{$V(h)$}
%      \makecell[l]{$r\leftarrow\Zq$\\$a\deq g^r$} & \send{a} & \\
%      & \recv{c} & $c\leftarrow\Zq$ \\
%      $z\deq r+c\cdot w$ & \send{z} & Output 1 if $g^z=a\cdot h^c$ \\
%    \end{sigmatab}
%    \end{protocolbox}
\newenvironment{sigmatab}[2]{%
  \small\renewcommand{\arraystretch}{1.6}%
  \begin{tabular}{@{}p{0.37\linewidth} c p{0.41\linewidth}@{}}%
    \multicolumn{1}{c}{\emph{Prover}\quad #1} & &
    \multicolumn{1}{c}{\emph{Verifier}\quad #2}\\[-3pt]
    \hline\\[-9pt]%
}{%
  \end{tabular}%
}

%  Attack / warning note box
\newenvironment{attacknote}{%
  \begin{tcolorbox}[
    enhanced,
    colback=red!5!white, colframe=red!55!black,
    boxrule=0.5pt, arc=2mm,
    left=5pt, right=5pt, top=3pt, bottom=3pt,
    fontupper=\small
  ]%
}{%
  \end{tcolorbox}%
}


\title{\textbf{Introduction to Zero Knowledge Proofs}}
\date{}

\begin{document}
\maketitle
\thispagestyle{empty}
\tableofcontents
\newpage

%% ============================================================
\section{What Is Zero Knowledge?}
%% ============================================================

In an ordinary proof you convince someone by showing them the evidence directly --- to prove you are over 18, you hand over your ID. The verifier learns your exact date of birth, your name, your address. The proof works, but it reveals far more than necessary.

Zero-knowledge proofs ask: can we do better? Can a prover convince a verifier that a statement is true while revealing \emph{nothing at all} beyond the bare truth of the claim?

The answer, perhaps surprisingly, is yes.

\subsection*{The simulator test}

Take \emph{perfect secrecy} as a concrete illustration. A cipher $\mathsf{Enc}$ is perfectly secret if
\[
  \Pr[m \mid c] = \Pr[m]
\]
for every message $m$ and every ciphertext $c$: observing the ciphertext does not shift your beliefs about the plaintext at all. The one-time pad is the canonical example --- the ciphertext is uniformly distributed regardless of the message, so it leaks nothing.

There is another way to see why this means no information leaks. Ask: could the adversary have produced the ciphertext themselves, without interacting with the encryption scheme at all? For the one-time pad, the answer is yes --- just sample a uniformly random string. The adversary can simulate the ciphertext on their own, so seeing the real one teaches them nothing new. This is the \emph{simulation paradigm}: if you can reproduce the output by yourself, the real output is useless to you.

\begin{exercisebox}
  Replace ``identical distributions'' with ``computationally indistinguishable distributions'' in the argument above. Write down the resulting security game between an adversary and a challenger, and verify that it captures CPA security for symmetric encryption.
\end{exercisebox}

\begin{tcolorbox}[colback=blue!4!white,colframe=blue!45!black]
  A proof system is \textbf{zero knowledge} if there exists an efficient \emph{simulator} that, \emph{without knowing the secret witness}, can produce transcripts that are indistinguishable from real interactions with the honest prover.
\end{tcolorbox}

The logic is elegant: if the verifier could have generated the same conversation by themselves (by running the simulator), then seeing the real conversation taught them nothing new. Whatever information seems to be in the proof was already accessible to the verifier without the prover's help.

We will develop this intuition through a series of concrete examples.

%% ============================================================
\section{The Ali Baba's Cave}
%% ============================================================

\subsection{Setup}

Imagine a circular cave with a single public entrance. Inside, the path splits into two branches, $A$ and $B$, that meet at a locked door deep in the rock. The door can only be opened with a secret magic word. Alice knows the magic word; Bob does not.

Alice wants to prove to Bob that she knows the magic word --- without uttering it.

\subsection{Protocol}

While Bob waits at the entrance, Alice enters the cave and walks down one of the two branches, choosing privately which one. Bob then shouts a challenge from the entrance, naming one of the two sides. Alice must emerge from that side: if she chose the correct branch she simply walks out; otherwise she must pass through the door using the magic word.

\medskip
\noindent
\begin{minipage}[t]{0.60\linewidth}
\begin{protocolbox}
\begin{sigmatab}{$P(\mathit{Cave};\ \mathit{secret})$}{$V(\mathit{Cave})$}
  \makecell[l]{Alice enters the cave and \\ chooses side $A$ or $B$ \\ (choice is private)} &
  \send{} & \\[3pt]
  & \recv{c} & $c \leftarrow \{A,\, B\}$ \\[3pt]
  \makecell[l]{Alice exits from side $z$ \\ (uses magic word if needed)} &
  \send{z} &
  \makecell[l]{Output 1 if $c = z$, \\ else output 0} \\
\end{sigmatab}
\end{protocolbox}
\end{minipage}\hfill
\begin{minipage}[t]{0.36\linewidth}
\vspace{0.5em}
\propbox{\propyes}{\propyes}{\propno}
\end{minipage}

\medskip

\subsection{Analysis}

\textbf{Completeness.} If Alice knows the magic word she can always exit from whichever side Bob names, regardless of which side she entered. The verifier always outputs 1.

\textbf{Honest-verifier zero knowledge.} The observable transcript is just $(c, z)$. An honest Alice always has $z = c$, so the transcript contains only the challenge value --- which the verifier chose themselves. A simulator that picks $c \leftarrow \{A, B\}$ at random and sets $z = c$ produces the same distribution without knowing any magic word.

\textbf{Special soundness fails.} Special soundness requires: given two accepting transcripts sharing the \emph{same commitment} but with \emph{different} challenges, an efficient extractor can recover the witness. Here the ``commitment'' is Alice's choice of branch, which is private and never observed by Bob. Even if we could somehow observe it, the transcripts only confirm that Alice exited the correct side --- they carry no algebraic trace of the magic word. No extractor can recover an opaque password from a pair of exit choices.

This is unlike the DLOG protocols below, where the witness $w$ appears \emph{linearly} in the response and can be recovered by subtraction.

%% ============================================================
\subsection{Exercise: The Apple Test}
%% ============================================================

\begin{exercisebox}
  Alice has two apples: one red, one green. Bob is red-green colorblind and cannot visually distinguish them. Alice claims she \emph{can} tell them apart. Describe a protocol by which Alice can convince Bob she has the ability to distinguish the apples, without revealing to Bob which apple is red and which is green.
\end{exercisebox}

%% ============================================================
\section{Sigma Protocols for Discrete Logarithm}
%% ============================================================

We now build a mathematically rigorous protocol for a number-theoretic statement, arriving at the famous \emph{Schnorr identification scheme}. Each attempt below fixes exactly one flaw of the previous one; the sequence of failures makes clear why every design choice in the final protocol is necessary.

\subsection{Setup}

Let $(\G, g, q)$ be a cyclic group of prime order $q$ with generator $g$. Write $\Z_q = \Z/q\Z$. The \emph{discrete logarithm (DLOG) relation} is
\[
  \relation = \bigl\{\,(h,\,w) \mid h = g^w\,\bigr\}.
\]
The prover $P$ holds a public element $h \in \G$ and a secret witness $w \in \Z_q$ with $h = g^w$. The verifier $V$ holds only $h$.

The goal: $P$ convinces $V$ that they know $w$ without revealing anything about $w$ itself.

%% -----------------------------------------------------------
\subsection{First Attempt: Reveal the Witness Directly}
%% -----------------------------------------------------------

The most direct approach: just send $w$.

\medskip
\noindent
\begin{minipage}[t]{0.60\linewidth}
\begin{protocolbox}
\begin{sigmatab}{$P(h;\, w)$}{$V(h)$}
  & \send{w} & Output 1 if $h = g^w$ \\
\end{sigmatab}
\end{protocolbox}
\end{minipage}\hfill
\begin{minipage}[t]{0.36\linewidth}
\vspace{0.5em}
\propbox{\propyes}{\propno}{\propyes}
\end{minipage}

\medskip

\textbf{Completeness} holds: the verifier checks $h \stackrel{?}{=} g^w$, true by assumption.

\textbf{Special soundness} holds trivially: the response $w$ is itself the witness.

\textbf{Honest-verifier ZK fails}: $w$ is sent in the clear. The verifier directly learns the discrete logarithm.

%% -----------------------------------------------------------
\subsection{Second Attempt: Blind the Response with a Nonce}
%% -----------------------------------------------------------

To hide $w$, the prover randomises the response with a fresh nonce $r$: commit to $g^r$ and then send $r$ as the response.

\medskip
\noindent
\begin{minipage}[t]{0.60\linewidth}
\begin{protocolbox}
\begin{sigmatab}{$P(h;\, w)$}{$V(h)$}
  \makecell[l]{$r \leftarrow \Zq$ \\ $a \deq g^r$} & \send{a} & \\[2pt]
  $z \deq r$ & \send{z} & Output 1 if $g^z = a$ \\
\end{sigmatab}
\end{protocolbox}
\end{minipage}\hfill
\begin{minipage}[t]{0.36\linewidth}
\vspace{0.5em}
\propbox{\propyes}{\propyes}{\propno}
\end{minipage}

\medskip

\textbf{Completeness} holds: $g^z = g^r = a$.

\textbf{Honest-verifier ZK} holds: $z = r$ is uniform in $\Z_q$ regardless of $w$. A simulator reproduces the transcript $(g^r, r)$ for any $r \leftarrow \Zq$ without knowing $w$.

\textbf{Soundness fails entirely.} The check $g^z \stackrel{?}{=} a$ does not involve $h$ at all. Anyone, regardless of whether they know $w$, can pick $r \leftarrow \Zq$, set $a = g^r$ and $z = r$, and the verifier accepts. This protocol proves nothing.

%% -----------------------------------------------------------
\subsection{Third Attempt: Incorporate the Witness into the Response}
%% -----------------------------------------------------------

The fix: include $w$ in the response via $z \deq r + w$ and adjust the check to involve $h$.

\medskip
\noindent
\begin{minipage}[t]{0.60\linewidth}
\begin{protocolbox}
\begin{sigmatab}{$P(h;\, w)$}{$V(h)$}
  \makecell[l]{$r \leftarrow \Zq$ \\ $a \deq g^r$} & \send{a} & \\[2pt]
  $z \deq r + w$ & \send{z} & Output 1 if $g^z = a \cdot h$ \\
\end{sigmatab}
\end{protocolbox}
\end{minipage}\hfill
\begin{minipage}[t]{0.36\linewidth}
\vspace{0.5em}
\propbox{\propyes}{\propyes}{\propno}
\end{minipage}

\medskip

\textbf{Completeness} holds: $g^z = g^{r+w} = g^r \cdot g^w = a \cdot h$.

\textbf{Honest-verifier ZK} holds: $z = r + w$ is uniform in $\Z_q$ since $r$ is uniform, so the transcript leaks no information about $w$.

\textbf{Soundness still fails.} A cheating prover can forge a convincing transcript \emph{without knowing $w$}:

\begin{attacknote}
\textbf{Attack.} Choose any $z \leftarrow \Zq$, then set $a \deq g^z \cdot h^{-1}$. Send $a$, then send $z$. The check becomes $g^z \stackrel{?}{=} a \cdot h = g^z \cdot h^{-1} \cdot h = g^z$. It passes.
\end{attacknote}

\subsection{What Is Missing?}

In Attempts 2 and 3 the prover sends both messages with \emph{no input from the verifier in between}. This freedom is fatal: the prover can choose the response $z$ first, then backwards-engineer the commitment $a$ to make the check pass --- bypassing the need to know $w$ entirely.

\begin{tcolorbox}[colback=blue!4!white,colframe=blue!45!black]
\textbf{Key insight.} The verifier must send a \emph{random challenge} between the commitment and the response. This forces the prover to fix $a$ before seeing the challenge, preventing any backwards engineering. A commitment chosen \emph{before} the challenge cannot be adapted to it.
\end{tcolorbox}

This is the defining structure of a \emph{sigma protocol}: commit $\to$ challenge $\to$ respond.

%% -----------------------------------------------------------
\subsection{Fourth Attempt: A Binary Challenge}
%% -----------------------------------------------------------

We combine the two verification checks from Attempts 2 and 3, selected by a random bit from the verifier.

\medskip
\noindent
\begin{minipage}[t]{0.65\linewidth}
\begin{protocolbox}
\begin{sigmatab}{$P(h;\, w)$}{$V(h)$}
  \makecell[l]{$r \leftarrow \Zq$ \\ $a \deq g^r$} & \send{a} & \\[2pt]
  & \recv{c} & $c \leftarrow \{0,\, 1\}$ \\[2pt]
  \makecell[l]{If $c=0$:\; $z \deq r$ \\ If $c=1$:\; $z \deq r+w$} &
  \send{z} &
  \makecell[l]{If $c=0$: output 1 if $g^z = a$ \\
               If $c=1$: output 1 if $g^z = a \cdot h$} \\
\end{sigmatab}
\end{protocolbox}
\end{minipage}\hfill
\begin{minipage}[t]{0.36\linewidth}
\vspace{0.5em}
\propbox{\propyes}{\propyes}{\propyes}
\vspace{0.5em}
{\small\color{orange!75!black} Soundness error: $1/2$.}
\end{minipage}

\medskip

\textbf{Completeness}: for $c=0$, $g^z = g^r = a$; for $c=1$, $g^z = g^{r+w} = a\cdot h$.

\textbf{Honest-verifier ZK}: for each fixed $c$, a simulator picks $z \leftarrow \Zq$, sets $a = g^z$ (if $c=0$) or $a = g^z \cdot h^{-1}$ (if $c=1$), and outputs $(a,c,z)$. Both distributions match the real transcript.

\begin{tcolorbox}[enhanced, colback=teal!5!white, colframe=teal!50!black, boxrule=0.6pt, arc=2mm, left=5pt, right=5pt, top=4pt, bottom=4pt]
Why is the simulator allowed to choose $a$ \emph{after} seeing $c$, while a real prover must commit to $a$ first?

\medskip
\textit{Reason 1 --- soundness.} If a simulator with no extra power could already fake transcripts, then any cheating prover could run that same strategy and break soundness. The simulator must therefore be given an advantage that a real prover does not have.

\medskip
\textit{Reason 2 --- who is the adversary.} Zero knowledge asks whether the \emph{verifier} can reproduce the transcript on their own. The verifier is the party trying to learn the witness, so it plays the role of the adversary in the simulation. An honest verifier samples $c$ themselves and therefore already knows $c$ before the transcript is complete. The simulator models exactly this honest verifier, and so is entitled to know $c$ in advance.
\end{tcolorbox}

\medskip

\begin{tcolorbox}[enhanced, colback=teal!5!white, colframe=teal!50!black, boxrule=0.6pt, arc=2mm, left=5pt, right=5pt, top=4pt, bottom=4pt]
This style of argument is called \emph{black-box simulation}: the simulator is given oracle access to the verifier. In the honest-verifier setting this is clean --- the verifier's randomness is fixed and used honestly, so the challenge is known to the simulator ahead of time. If the verifier is malicious, it may use its randomness arbitrarily, and the simulator can no longer rely on knowing $c$ in advance. Handling a malicious verifier requires a more powerful technique called \emph{rewinding}, which is beyond our scope here.
\end{tcolorbox}

\textbf{Special soundness}: given two accepting transcripts $(a, 0, z_0)$ and $(a, 1, z_1)$ with the same $a$, we have $z_0 = r$ and $z_1 = r + w$. Subtracting:
\[
  w = z_1 - z_0 \pmod{q}.
\]

A cheating prover without $w$ can still succeed with probability $1/2$: they commit to an $a$ valid for one challenge value and hope the verifier asks for it. Running $k$ independent rounds reduces this to $2^{-k}$, but there is a cleaner way.

%% -----------------------------------------------------------
\subsection{The Schnorr Protocol: Large Challenge Space}
%% -----------------------------------------------------------

Replace the binary challenge with $c \leftarrow \Zq$. The soundness error drops from $1/2$ to $1/q$, negligible for a cryptographically large prime $q$.

\medskip
\noindent
\begin{minipage}[t]{0.60\linewidth}
\begin{protocolbox}
\begin{sigmatab}{$P(h;\, w)$}{$V(h)$}
  \makecell[l]{$r \leftarrow \Zq$ \\ $a \deq g^r$} & \send{a} & \\[2pt]
  & \recv{c} & $c \leftarrow \Zq$ \\[2pt]
  $z \deq r + c \cdot w$ & \send{z} &
  Output 1 if $g^z = a \cdot h^c$ \\
\end{sigmatab}
\end{protocolbox}
\end{minipage}\hfill
\begin{minipage}[t]{0.36\linewidth}
\vspace{0.5em}
\propbox{\propyes}{\propyes}{\propyes}
\vspace{0.5em}
{\small\color{green!60!black} Soundness error: $1/q$ (negligible).}
\end{minipage}

\medskip

\textbf{Completeness}:
\[
  g^z = g^{r+cw} = g^r \cdot (g^w)^c = a \cdot h^c. \qquad\checkmark
\]

\textbf{Honest-verifier ZK}: for any fixed $c$, the simulator picks $z \leftarrow \Zq$, sets $a \deq g^z \cdot h^{-c}$, and outputs $(a, c, z)$. Check: $g^z \stackrel{?}{=} a \cdot h^c = g^z$. The joint distribution of $(a,z)$ is uniform over $\G \times \Zq$ --- identical to the real protocol.

\textbf{Special soundness}: given two accepting transcripts $(a, c, z)$ and $(a, c', z')$ with $c \neq c'$:
\[
  z = r + cw, \qquad z' = r + c'w
  \qquad\Longrightarrow\qquad
  w = \frac{z - z'}{c - c'} \pmod{q}.
\]
Division is well-defined because $q$ is prime and $c \neq c'$ implies $c - c'$ is invertible.

\medskip

\begin{tcolorbox}[colback=green!4!white,colframe=green!50!black,boxrule=0.6pt]
\textbf{Identification scheme.} The Schnorr protocol serves as a public-key identification scheme. The prover's \emph{secret key} is $w \leftarrow \Zq$; their \emph{public key} is $h = g^w$. Since computing $w$ from $h$ requires solving the discrete logarithm --- hard in suitable groups --- only the party who generated $h$ knows $w$. Each run of the protocol lets the prover authenticate as the owner of $h$ without exposing any information about $w$.
\end{tcolorbox}

%% ============================================================
\section{A Fundamental Tension: Soundness vs.\ Zero Knowledge}
%% ============================================================

Before moving to the non-interactive world it is worth pausing on a basic incompatibility.

\begin{definition}[Perfect soundness]
A proof system has \emph{perfect soundness} if no cheating prover --- even one with unbounded computation --- can convince the verifier of a false statement with any nonzero probability.
\end{definition}

\begin{definition}[Zero knowledge]
A proof system is \emph{zero knowledge} if there exists an efficient simulator that, without the witness, produces transcripts that are computationally indistinguishable from real ones.
\end{definition}

\textbf{Why they cannot coexist.} The core mechanism behind zero knowledge is simple: if the simulator knows the challenge $c$ \emph{before} it commits to $a$, it can always produce an indistinguishable transcript without knowing the witness. In Schnorr, the simulator picks $z \leftarrow \Zq$ freely and back-computes $a \deq g^z \cdot h^{-c}$. Because $z$ is uniform and the computation is purely algebraic, the resulting transcript is indistinguishable from a real one (here it is even identically distributed).

Now contrast this with perfect soundness: if the proof system is perfectly sound, then for any false statement $h \notin \lang$, even a prover who knows the challenge $c$ in advance should have \emph{zero} probability of producing an accepting transcript. But the simulator technique above always succeeds for any $h$ as long as the challenge $c$ is known beforehand, without knowing the witness. The two requirements are directly at odds.

The resolution is to allow a small soundness error. Think of it as a \emph{backdoor}: the same algebraic shortcut that lets the simulator bypass the witness also lets a cheating prover occasionally sneak through. In Schnorr, a cheating prover can guess the challenge $c$ correctly with probability $1/q$ before committing, and then apply the simulator trick to produce an accepting transcript. That $1/q$ leakage is the unavoidable price of zero knowledge.

%% ============================================================
\section{From Interactive to Non-Interactive: The Role of the CRS}
%% ============================================================

\subsection{The problem with NIZK}

In the sigma protocols above, soundness rests on the random challenge $c$ sent by the verifier \emph{after} the prover commits. Recall Attempts 2 and 3: without that challenge, the prover could choose the response first and backwards-engineer a matching commitment. The challenge is the essential ingredient that prevents this.

In a \emph{non-interactive zero-knowledge} (NIZK) proof system there is no back-and-forth. The prover sends a single proof string $\zkproof$; the verifier checks it with no further communication. There is no external $c$ to anchor the commitment.

\begin{tcolorbox}[colback=red!5!white,colframe=red!55!black,boxrule=0.6pt,arc=2mm]
\textbf{The difficulty.} Without a verifier challenge, the prover and the ZK simulator face exactly the same task: produce $\zkproof$ without receiving anything from the outside. If a simulator can generate convincing proofs for true statements without the witness, a cheating prover can attempt the same strategy for false statements. Soundness and zero knowledge appear to be incompatible.
\end{tcolorbox}

\subsection{The common reference string}

The solution is to introduce shared randomness via a \emph{common reference string} (CRS). A trusted setup algorithm generates
\[
  (\crs,\, \td) \leftarrow \simsetup(1^\lambda,\, \relation),
\]
publishes $\crs$, and discards (or keeps secret) the \emph{trapdoor} $\td$. The crucial asymmetry is:

\begin{itemize}[leftmargin=2em]
  \item A real prover sees only $\crs$ and must know the witness $w$ to produce a valid proof.
  \item The ZK simulator additionally knows the trapdoor $\td$, which lets it produce fake proofs even without $w$ --- or even for false statements.
\end{itemize}

This restores the separation that the interactive challenge provided. In the interactive setting, the simulator's extra power was the ability to see $c$ before committing. In the NIZK setting, the simulator's extra power is the trapdoor $\td$.

%% groth16.tex -- Construction of the Groth16 SNARK
%% %% groth16.tex -- Construction of the Groth16 SNARK
%% %% groth16.tex -- Construction of the Groth16 SNARK
%% \input{groth16} inside zero_knowledge.tex

%% ============================================================
\section{Groth16: A Pairing-Based SNARK}
%% ============================================================

Groth16 is a preprocessing SNARK for arithmetic circuit satisfiability. Its proof is only three group elements, verification is a single pairing equation, and it achieves perfect completeness and perfect zero knowledge. Soundness holds against adversaries that interact with the system via generic group operations.

We build the system in three layers: encode computation as a polynomial divisibility problem (Quadratic Arithmetic Program, QAP), construct an information-theoretic proof for it (Non-Interactive Linear Proof, NILP), then compile the NILP into group elements using pairings.

%% -----------------------------------------------------------
\subsection{Notation for Pairing Groups}
%% -----------------------------------------------------------

Let $(p, \G_1, \G_2, \GT, e, g, h)$ be an asymmetric bilinear group of prime order $p$. The generator of $\G_1$ is $g$, the generator of $\G_2$ is $h$, and $e : \G_1 \times \G_2 \to \GT$ is a bilinear pairing satisfying $e(g^a, h^b) = e(g,h)^{ab}$.

We represent group elements by their discrete logarithms using the bracket notation
\[
  \gone{a} \deq g^a \in \G_1, \qquad
  \gtwo{a} \deq h^a \in \G_2, \qquad
  \gT{a}   \deq e(g,h)^a \in \GT, \qquad a \in \Z_p.
\]
This notation makes linear algebra readable. For a known scalar $c$ or known matrix $M$:
\[
  \gone{a} + \gone{b} = \gone{a+b}, \qquad
  c \cdot \gone{a} = \gone{ca}, \qquad
  M \cdot \gone{\mathbf{a}} = \gone{M\mathbf{a}}.
\]
The pairing provides one level of multiplication across the two source groups:
\[
  \gone{a} \cdot \gtwo{b} \deq e\!\left(g^a, h^b\right) = \gT{ab}.
\]


%% -----------------------------------------------------------
\subsection{From Arithmetic Circuits to Quadratic Arithmetic Programs}
%% -----------------------------------------------------------

\subsubsection*{Arithmetic circuits and R1CS}

An arithmetic circuit over a finite field $\F_p$ consists of addition and multiplication gates over field elements. Any such circuit can be \emph{flattened}: every gate is rewritten as a single multiplicative constraint
\[
  \bigl(\textstyle\sum_i A_{q,i}\, a_i\bigr)
  \cdot
  \bigl(\textstyle\sum_i B_{q,i}\, a_i\bigr)
  =
  \textstyle\sum_i C_{q,i}\, a_i
  \qquad (q = 1,\ldots, n),
\]
where $z = (a_0, a_1, \ldots, a_m) \in \F_p^{m+1}$ is the \emph{assignment vector}: $a_0 = 1$ is a constant, $\phi = (a_1, \ldots, a_\ell)$ is the public statement, and $w = (a_{\ell+1}, \ldots, a_m)$ is the secret witness. The $n$ constraints are collected into three matrices $A, B, C \in \F_p^{n \times (m+1)}$ --- this is the \emph{rank-1 constraint system} (R1CS). Arithmetic circuit satisfiability is NP-complete: any NP statement can be compiled into an R1CS instance over a suitable field, so a proof system for R1CS is a proof system for all of NP.

\subsubsection*{Quadratic arithmetic programs}

Checking $n$ separate scalar equations one by one would be costly. The key algebraic trick is to compress them into a single polynomial divisibility test.

Fix $n$ distinct field elements $r_1, \ldots, r_n \in \F_p$, one for each constraint. Define the vanishing polynomial
\[
  t(X) = \prod_{q=1}^{n}(X - r_q).
\]
For each variable index $i \in \{0, \ldots, m\}$, interpolate three polynomials $u_i, v_i, w_i \in \F_p[X]$ of degree at most $n-1$ such that
\[
  u_i(r_q) = A_{q,i}, \qquad v_i(r_q) = B_{q,i}, \qquad w_i(r_q) = C_{q,i} \qquad \forall\, q.
\]
\todo{Add a visualisation showing how the columns of the $A$, $B$, $C$ matrices are read off and interpolated into the polynomials $u_i$, $v_i$, $w_i$.}
Given an assignment $z$, define the aggregated polynomials
\[
  U(X) = \sum_{i=0}^{m} a_i\, u_i(X), \qquad
  V(X) = \sum_{i=0}^{m} a_i\, v_i(X), \qquad
  W(X) = \sum_{i=0}^{m} a_i\, w_i(X).
\]
At each constraint point $r_q$, the R1CS equation $U(r_q)\cdot V(r_q) = W(r_q)$ must hold. Since $t(X)$ vanishes precisely at $r_1, \ldots, r_n$, this is equivalent to
\[
  \boxed{t(X) \;\Big|\; U(X)\cdot V(X) - W(X).}
\]
Equivalently, there exists a degree-$(n-2)$ quotient polynomial $h(X)$ such that
\[
  U(X)\cdot V(X) = W(X) + h(X)\,t(X).
\]
This is the \emph{quadratic arithmetic program} (QAP) relation for the circuit. The advantage over R1CS is that $n$ separate equations have been replaced by one polynomial identity checkable at a single random point --- exactly the Schwartz--Zippel trick.

%% -----------------------------------------------------------
\subsection{A Non-Interactive Linear Proof for QAP}
%% -----------------------------------------------------------

A \emph{non-interactive linear proof} (NILP) is an information-theoretic proof system where the prover's messages are linear functions of a common reference string $\sigma$. We now construct such a system for the QAP relation.

\todo{Explain why this NILP is non-interactive despite stemming from a Linear Interactive Proof (LIP): the CRS plays the role of the verifier's challenges, collapsing the interaction into a single message. Discuss the connection between 2-move LIPs and NILPs more carefully.}

\subsubsection*{Setup}

Pick a random trapdoor $\tau = (\alpha, \beta, \gamma, \delta, x) \in \F_p^{*5}$ and define the common reference string $\sigma$ as the collection of field evaluations
\[
\sigma =
\left(
\begin{aligned}
&\alpha,\; \beta,\; \gamma,\; \delta,\;\{x^i\}_{i=0}^{n-1},\\[2pt]
&\left\{\frac{\beta u_i(x)+\alpha v_i(x)+w_i(x)}{\gamma}\right\}_{i=0}^{\ell},\\[4pt]
&\left\{\frac{\beta u_i(x)+\alpha v_i(x)+w_i(x)}{\delta}\right\}_{i=\ell+1}^{m},\\[4pt]
&\left\{\frac{x^i\, t(x)}{\delta}\right\}_{i=0}^{n-2}
\end{aligned}
\right).
\]
The first block enables evaluation of arbitrary polynomials at $x$. The second block encodes the public witness terms, scaled by $\gamma$. The third encodes the private witness terms, scaled by $\delta$. The last encodes powers of $x$ weighted by $t(x)/\delta$, allowing the prover to compute $h(x)\,t(x)/\delta$ as a linear combination.

\todo{Discuss what $\alpha$, $\beta$, $\gamma$, and $\delta$ each do exactly: why they appear in the CRS, how they enforce consistency between $A$, $B$, $C$, and why the $\gamma$/$\delta$ separation is necessary.}

\subsubsection*{Prover}

Given a satisfying assignment $(a_1, \ldots, a_m)$, the prover computes $h(X)$ from the QAP relation, then picks two random blinding scalars $r, s \leftarrow \F_p$ and sets
\begin{align*}
  A &\deq \alpha + U(x) + r\delta, \\
  B &\deq \beta  + V(x) + s\delta, \\
  C &\deq \frac{\displaystyle\sum_{i=\ell+1}^{m} a_i\bigl(\beta u_i(x)+\alpha v_i(x)+w_i(x)\bigr) + h(x)\,t(x)}{\delta} + As + rB - rs\delta.
\end{align*}
The proof is $\pi = (A, B, C)$.

\subsubsection*{Verifier}

The verifier checks the single equation
\[
  A \cdot B
  \;\stackrel{?}{=}\;
  \underbrace{\alpha\beta}_{\text{(i)}}
  \;+\;
  \underbrace{\sum_{i=0}^{\ell} a_i \cdot \frac{\beta u_i(x)+\alpha v_i(x)+w_i(x)}{\gamma}}_{\text{(ii)}} \cdot\, \gamma
  \;+\;
  \underbrace{C \cdot \delta}_{\text{(iii)}}.
\]
The three terms on the right play distinct roles.
\begin{itemize}[noitemsep, leftmargin=2em]
  \item[(i)] $\alpha\beta$ is a fixed anchor from the setup. Its presence forces $A$ and $B$ to contain genuine $\alpha$ and $\beta$ offsets respectively, preventing a prover from constructing $A$ and $B$ independently of the structure of $\sigma$.
  \item[(ii)] The sum over $i = 0, \ldots, \ell$ incorporates the \emph{public inputs} $a_0, \ldots, a_\ell$ --- the statement $\phi$ that the verifier knows. These are precomputed in the CRS scaled by $\gamma$ and $\gamma$ cancels out, leaving the verifier to evaluate this term using only public data. Note that $\alpha$ and $\beta$ also appear here, but they are \emph{setup parameters} embedded in the precomputed CRS entries, not the circuit's inputs.
  \item[(iii)] $C \cdot \delta$ is where the private witness $a_{\ell+1}, \ldots, a_m$ is encoded. The $\delta$ scaling keeps it algebraically separated from term (ii), preventing the prover from mixing public and private contributions.
\end{itemize}

\subsubsection*{Simulator}

Given the trapdoor $\tau$, the simulator ignores the witness entirely. It picks $A, B \leftarrow \F_p$ uniformly at random and computes
\[
  C \deq \frac{A B \;-\; \alpha\beta \;-\; \displaystyle\sum_{i=0}^{\ell} a_i\bigl(\beta u_i(x)+\alpha v_i(x)+w_i(x)\bigr)}{\delta}.
\]

\subsubsection*{Why this achieves zero knowledge}

The blinding scalars $r$ and $s$ make $A$ and $B$ uniformly distributed in $\F_p$, independent of the witness. Given uniform $A$ and $B$, the verification equation uniquely determines $C$. The real distribution of $(A, B, C)$ is therefore identical to the simulated one: two uniform field elements followed by their forced companion. This is perfect zero knowledge.

\subsubsection*{Why this is sound against affine strategies}

An affine prover strategy computes $(A, B, C)$ as linear combinations of $\sigma$. Viewing the verification equation as a polynomial identity in the formal indeterminates $\alpha, \beta, \gamma, \delta, x$ and applying Schwartz--Zippel, any affine strategy that passes verification with non-negligible probability must satisfy the identity formally. A careful analysis of the monomials (particularly those involving $\alpha\beta$, $1/\gamma^2$, and $1/\delta^2$) forces the coefficients to encode a consistent assignment, from which the witness $(a_{\ell+1}, \ldots, a_m)$ can be extracted.

\todo{Add the full monomial analysis proof here, following Theorem~1 of the Groth (2016) paper.}

%% -----------------------------------------------------------
\subsection{Compiling to Pairings: The Generic Group Model}
%% -----------------------------------------------------------

The NILP works over a field, but we need a cryptographic argument --- the prover should not be able to read off $x$ from $\sigma$. The compilation step replaces every field element in $\sigma$ with its bracket encoding, and $(A, B, C)$ become $(\gone{A}, \gtwo{B}, \gone{C})$.

The compiled verification equation is:
\[
  \gone{A} \cdot \gtwo{B}
  \;=\;
  \gone{\alpha}\cdot\gtwo{\beta}
  \;+\;
  \sum_{i=0}^{\ell} a_i\,
    \left[\frac{\beta u_i(x)+\alpha v_i(x)+w_i(x)}{\gamma}\right]_{\!1}
    \cdot\, \gtwo{\gamma}
  \;+\;
  \gone{C}\cdot\gtwo{\delta}.
\]
The prover computes $\gone{A}$, $\gone{C}$, $\gtwo{B}$ as linear combinations of the CRS elements --- this is exactly the NILP prover running in the exponent. The verifier checks one pairing product equation using three pairings in total.

\textbf{Lifting soundness via the generic group model.} In the generic group model, a prover can only use the group elements in $\sigma$ through addition (and scalar multiplication) --- it cannot read their discrete logarithms. Any computation it performs on $\gone{\sigma_1}$ and $\gtwo{\sigma_2}$ is therefore a sequence of affine operations, corresponding exactly to choosing a proof matrix $\Pi$ and computing $\Pi\sigma$. This is precisely an affine strategy in the NILP sense. Combined with a \emph{disclosure-freeness} property of the CRS --- which guarantees that the structure of $\sigma$ cannot be exploited to choose a non-affine strategy --- the NILP's soundness against affine strategies lifts to soundness against all generic adversaries.

%% -----------------------------------------------------------
\subsection*{Summary: The Groth16 Proof System}
%% -----------------------------------------------------------

\begin{tcolorbox}[enhanced, colback=blue!3!white, colframe=blue!45!black,
                  arc=2.5pt, boxrule=1pt, left=6pt, right=6pt, top=5pt, bottom=5pt]
\small
\textbf{Setup.} Sample $\tau = (\alpha,\beta,\gamma,\delta,x) \leftarrow \F_p^{*5}$.
Publish $\sigma = \bigl([\sigma_1]_1,\, [\sigma_2]_2\bigr)$ as described above. Keep $\tau$ secret.

\medskip
\textbf{Prove.} Given $(a_1,\ldots,a_m)$ satisfying the circuit, pick $r,s \leftarrow \F_p$ and output
$\pi = \bigl(\gone{A},\, \gone{C},\, \gtwo{B}\bigr) \in \G_1^2 \times \G_2$.

\medskip
\textbf{Verify.} Check:
$\gone{A} \cdot \gtwo{B} = \gone{\alpha}\cdot\gtwo{\beta} + \displaystyle\sum_{i=0}^{\ell} a_i\,\gone{(\cdot)_i}\cdot\gtwo{\gamma} + \gone{C}\cdot\gtwo{\delta}$.

\medskip
\textbf{Proof size:} \quad $2$ elements in $\G_1$, \quad $1$ element in $\G_2$ \quad $= 3$ group elements.\\
\textbf{Verification:} \quad $1$ pairing product equation requiring $3$ online pairings:
$e(\gone{A}, \gtwo{B})$,\; $e\!\left(\textstyle\sum a_i K_i,\, \gtwo{\gamma}\right)$,\; $e(\gone{C}, \gtwo{\delta})$.
The fourth term $\gT{\alpha\beta}$ is precomputed and stored in the verifier key, so no extra pairing is needed for it.
\end{tcolorbox}
 inside zero_knowledge.tex

%% ============================================================
\section{Groth16: A Pairing-Based SNARK}
%% ============================================================

Groth16 is a preprocessing SNARK for arithmetic circuit satisfiability. Its proof is only three group elements, verification is a single pairing equation, and it achieves perfect completeness and perfect zero knowledge. Soundness holds against adversaries that interact with the system via generic group operations.

We build the system in three layers: encode computation as a polynomial divisibility problem (Quadratic Arithmetic Program, QAP), construct an information-theoretic proof for it (Non-Interactive Linear Proof, NILP), then compile the NILP into group elements using pairings.

%% -----------------------------------------------------------
\subsection{Notation for Pairing Groups}
%% -----------------------------------------------------------

Let $(p, \G_1, \G_2, \GT, e, g, h)$ be an asymmetric bilinear group of prime order $p$. The generator of $\G_1$ is $g$, the generator of $\G_2$ is $h$, and $e : \G_1 \times \G_2 \to \GT$ is a bilinear pairing satisfying $e(g^a, h^b) = e(g,h)^{ab}$.

We represent group elements by their discrete logarithms using the bracket notation
\[
  \gone{a} \deq g^a \in \G_1, \qquad
  \gtwo{a} \deq h^a \in \G_2, \qquad
  \gT{a}   \deq e(g,h)^a \in \GT, \qquad a \in \Z_p.
\]
This notation makes linear algebra readable. For a known scalar $c$ or known matrix $M$:
\[
  \gone{a} + \gone{b} = \gone{a+b}, \qquad
  c \cdot \gone{a} = \gone{ca}, \qquad
  M \cdot \gone{\mathbf{a}} = \gone{M\mathbf{a}}.
\]
The pairing provides one level of multiplication across the two source groups:
\[
  \gone{a} \cdot \gtwo{b} \deq e\!\left(g^a, h^b\right) = \gT{ab}.
\]


%% -----------------------------------------------------------
\subsection{From Arithmetic Circuits to Quadratic Arithmetic Programs}
%% -----------------------------------------------------------

\subsubsection*{Arithmetic circuits and R1CS}

An arithmetic circuit over a finite field $\F_p$ consists of addition and multiplication gates over field elements. Any such circuit can be \emph{flattened}: every gate is rewritten as a single multiplicative constraint
\[
  \bigl(\textstyle\sum_i A_{q,i}\, a_i\bigr)
  \cdot
  \bigl(\textstyle\sum_i B_{q,i}\, a_i\bigr)
  =
  \textstyle\sum_i C_{q,i}\, a_i
  \qquad (q = 1,\ldots, n),
\]
where $z = (a_0, a_1, \ldots, a_m) \in \F_p^{m+1}$ is the \emph{assignment vector}: $a_0 = 1$ is a constant, $\phi = (a_1, \ldots, a_\ell)$ is the public statement, and $w = (a_{\ell+1}, \ldots, a_m)$ is the secret witness. The $n$ constraints are collected into three matrices $A, B, C \in \F_p^{n \times (m+1)}$ --- this is the \emph{rank-1 constraint system} (R1CS). Arithmetic circuit satisfiability is NP-complete: any NP statement can be compiled into an R1CS instance over a suitable field, so a proof system for R1CS is a proof system for all of NP.

\subsubsection*{Quadratic arithmetic programs}

Checking $n$ separate scalar equations one by one would be costly. The key algebraic trick is to compress them into a single polynomial divisibility test.

Fix $n$ distinct field elements $r_1, \ldots, r_n \in \F_p$, one for each constraint. Define the vanishing polynomial
\[
  t(X) = \prod_{q=1}^{n}(X - r_q).
\]
For each variable index $i \in \{0, \ldots, m\}$, interpolate three polynomials $u_i, v_i, w_i \in \F_p[X]$ of degree at most $n-1$ such that
\[
  u_i(r_q) = A_{q,i}, \qquad v_i(r_q) = B_{q,i}, \qquad w_i(r_q) = C_{q,i} \qquad \forall\, q.
\]
\todo{Add a visualisation showing how the columns of the $A$, $B$, $C$ matrices are read off and interpolated into the polynomials $u_i$, $v_i$, $w_i$.}
Given an assignment $z$, define the aggregated polynomials
\[
  U(X) = \sum_{i=0}^{m} a_i\, u_i(X), \qquad
  V(X) = \sum_{i=0}^{m} a_i\, v_i(X), \qquad
  W(X) = \sum_{i=0}^{m} a_i\, w_i(X).
\]
At each constraint point $r_q$, the R1CS equation $U(r_q)\cdot V(r_q) = W(r_q)$ must hold. Since $t(X)$ vanishes precisely at $r_1, \ldots, r_n$, this is equivalent to
\[
  \boxed{t(X) \;\Big|\; U(X)\cdot V(X) - W(X).}
\]
Equivalently, there exists a degree-$(n-2)$ quotient polynomial $h(X)$ such that
\[
  U(X)\cdot V(X) = W(X) + h(X)\,t(X).
\]
This is the \emph{quadratic arithmetic program} (QAP) relation for the circuit. The advantage over R1CS is that $n$ separate equations have been replaced by one polynomial identity checkable at a single random point --- exactly the Schwartz--Zippel trick.

%% -----------------------------------------------------------
\subsection{A Non-Interactive Linear Proof for QAP}
%% -----------------------------------------------------------

A \emph{non-interactive linear proof} (NILP) is an information-theoretic proof system where the prover's messages are linear functions of a common reference string $\sigma$. We now construct such a system for the QAP relation.

\todo{Explain why this NILP is non-interactive despite stemming from a Linear Interactive Proof (LIP): the CRS plays the role of the verifier's challenges, collapsing the interaction into a single message. Discuss the connection between 2-move LIPs and NILPs more carefully.}

\subsubsection*{Setup}

Pick a random trapdoor $\tau = (\alpha, \beta, \gamma, \delta, x) \in \F_p^{*5}$ and define the common reference string $\sigma$ as the collection of field evaluations
\[
\sigma =
\left(
\begin{aligned}
&\alpha,\; \beta,\; \gamma,\; \delta,\;\{x^i\}_{i=0}^{n-1},\\[2pt]
&\left\{\frac{\beta u_i(x)+\alpha v_i(x)+w_i(x)}{\gamma}\right\}_{i=0}^{\ell},\\[4pt]
&\left\{\frac{\beta u_i(x)+\alpha v_i(x)+w_i(x)}{\delta}\right\}_{i=\ell+1}^{m},\\[4pt]
&\left\{\frac{x^i\, t(x)}{\delta}\right\}_{i=0}^{n-2}
\end{aligned}
\right).
\]
The first block enables evaluation of arbitrary polynomials at $x$. The second block encodes the public witness terms, scaled by $\gamma$. The third encodes the private witness terms, scaled by $\delta$. The last encodes powers of $x$ weighted by $t(x)/\delta$, allowing the prover to compute $h(x)\,t(x)/\delta$ as a linear combination.

\todo{Discuss what $\alpha$, $\beta$, $\gamma$, and $\delta$ each do exactly: why they appear in the CRS, how they enforce consistency between $A$, $B$, $C$, and why the $\gamma$/$\delta$ separation is necessary.}

\subsubsection*{Prover}

Given a satisfying assignment $(a_1, \ldots, a_m)$, the prover computes $h(X)$ from the QAP relation, then picks two random blinding scalars $r, s \leftarrow \F_p$ and sets
\begin{align*}
  A &\deq \alpha + U(x) + r\delta, \\
  B &\deq \beta  + V(x) + s\delta, \\
  C &\deq \frac{\displaystyle\sum_{i=\ell+1}^{m} a_i\bigl(\beta u_i(x)+\alpha v_i(x)+w_i(x)\bigr) + h(x)\,t(x)}{\delta} + As + rB - rs\delta.
\end{align*}
The proof is $\pi = (A, B, C)$.

\subsubsection*{Verifier}

The verifier checks the single equation
\[
  A \cdot B
  \;\stackrel{?}{=}\;
  \underbrace{\alpha\beta}_{\text{(i)}}
  \;+\;
  \underbrace{\sum_{i=0}^{\ell} a_i \cdot \frac{\beta u_i(x)+\alpha v_i(x)+w_i(x)}{\gamma}}_{\text{(ii)}} \cdot\, \gamma
  \;+\;
  \underbrace{C \cdot \delta}_{\text{(iii)}}.
\]
The three terms on the right play distinct roles.
\begin{itemize}[noitemsep, leftmargin=2em]
  \item[(i)] $\alpha\beta$ is a fixed anchor from the setup. Its presence forces $A$ and $B$ to contain genuine $\alpha$ and $\beta$ offsets respectively, preventing a prover from constructing $A$ and $B$ independently of the structure of $\sigma$.
  \item[(ii)] The sum over $i = 0, \ldots, \ell$ incorporates the \emph{public inputs} $a_0, \ldots, a_\ell$ --- the statement $\phi$ that the verifier knows. These are precomputed in the CRS scaled by $\gamma$ and $\gamma$ cancels out, leaving the verifier to evaluate this term using only public data. Note that $\alpha$ and $\beta$ also appear here, but they are \emph{setup parameters} embedded in the precomputed CRS entries, not the circuit's inputs.
  \item[(iii)] $C \cdot \delta$ is where the private witness $a_{\ell+1}, \ldots, a_m$ is encoded. The $\delta$ scaling keeps it algebraically separated from term (ii), preventing the prover from mixing public and private contributions.
\end{itemize}

\subsubsection*{Simulator}

Given the trapdoor $\tau$, the simulator ignores the witness entirely. It picks $A, B \leftarrow \F_p$ uniformly at random and computes
\[
  C \deq \frac{A B \;-\; \alpha\beta \;-\; \displaystyle\sum_{i=0}^{\ell} a_i\bigl(\beta u_i(x)+\alpha v_i(x)+w_i(x)\bigr)}{\delta}.
\]

\subsubsection*{Why this achieves zero knowledge}

The blinding scalars $r$ and $s$ make $A$ and $B$ uniformly distributed in $\F_p$, independent of the witness. Given uniform $A$ and $B$, the verification equation uniquely determines $C$. The real distribution of $(A, B, C)$ is therefore identical to the simulated one: two uniform field elements followed by their forced companion. This is perfect zero knowledge.

\subsubsection*{Why this is sound against affine strategies}

An affine prover strategy computes $(A, B, C)$ as linear combinations of $\sigma$. Viewing the verification equation as a polynomial identity in the formal indeterminates $\alpha, \beta, \gamma, \delta, x$ and applying Schwartz--Zippel, any affine strategy that passes verification with non-negligible probability must satisfy the identity formally. A careful analysis of the monomials (particularly those involving $\alpha\beta$, $1/\gamma^2$, and $1/\delta^2$) forces the coefficients to encode a consistent assignment, from which the witness $(a_{\ell+1}, \ldots, a_m)$ can be extracted.

\todo{Add the full monomial analysis proof here, following Theorem~1 of the Groth (2016) paper.}

%% -----------------------------------------------------------
\subsection{Compiling to Pairings: The Generic Group Model}
%% -----------------------------------------------------------

The NILP works over a field, but we need a cryptographic argument --- the prover should not be able to read off $x$ from $\sigma$. The compilation step replaces every field element in $\sigma$ with its bracket encoding, and $(A, B, C)$ become $(\gone{A}, \gtwo{B}, \gone{C})$.

The compiled verification equation is:
\[
  \gone{A} \cdot \gtwo{B}
  \;=\;
  \gone{\alpha}\cdot\gtwo{\beta}
  \;+\;
  \sum_{i=0}^{\ell} a_i\,
    \left[\frac{\beta u_i(x)+\alpha v_i(x)+w_i(x)}{\gamma}\right]_{\!1}
    \cdot\, \gtwo{\gamma}
  \;+\;
  \gone{C}\cdot\gtwo{\delta}.
\]
The prover computes $\gone{A}$, $\gone{C}$, $\gtwo{B}$ as linear combinations of the CRS elements --- this is exactly the NILP prover running in the exponent. The verifier checks one pairing product equation using three pairings in total.

\textbf{Lifting soundness via the generic group model.} In the generic group model, a prover can only use the group elements in $\sigma$ through addition (and scalar multiplication) --- it cannot read their discrete logarithms. Any computation it performs on $\gone{\sigma_1}$ and $\gtwo{\sigma_2}$ is therefore a sequence of affine operations, corresponding exactly to choosing a proof matrix $\Pi$ and computing $\Pi\sigma$. This is precisely an affine strategy in the NILP sense. Combined with a \emph{disclosure-freeness} property of the CRS --- which guarantees that the structure of $\sigma$ cannot be exploited to choose a non-affine strategy --- the NILP's soundness against affine strategies lifts to soundness against all generic adversaries.

%% -----------------------------------------------------------
\subsection*{Summary: The Groth16 Proof System}
%% -----------------------------------------------------------

\begin{tcolorbox}[enhanced, colback=blue!3!white, colframe=blue!45!black,
                  arc=2.5pt, boxrule=1pt, left=6pt, right=6pt, top=5pt, bottom=5pt]
\small
\textbf{Setup.} Sample $\tau = (\alpha,\beta,\gamma,\delta,x) \leftarrow \F_p^{*5}$.
Publish $\sigma = \bigl([\sigma_1]_1,\, [\sigma_2]_2\bigr)$ as described above. Keep $\tau$ secret.

\medskip
\textbf{Prove.} Given $(a_1,\ldots,a_m)$ satisfying the circuit, pick $r,s \leftarrow \F_p$ and output
$\pi = \bigl(\gone{A},\, \gone{C},\, \gtwo{B}\bigr) \in \G_1^2 \times \G_2$.

\medskip
\textbf{Verify.} Check:
$\gone{A} \cdot \gtwo{B} = \gone{\alpha}\cdot\gtwo{\beta} + \displaystyle\sum_{i=0}^{\ell} a_i\,\gone{(\cdot)_i}\cdot\gtwo{\gamma} + \gone{C}\cdot\gtwo{\delta}$.

\medskip
\textbf{Proof size:} \quad $2$ elements in $\G_1$, \quad $1$ element in $\G_2$ \quad $= 3$ group elements.\\
\textbf{Verification:} \quad $1$ pairing product equation requiring $3$ online pairings:
$e(\gone{A}, \gtwo{B})$,\; $e\!\left(\textstyle\sum a_i K_i,\, \gtwo{\gamma}\right)$,\; $e(\gone{C}, \gtwo{\delta})$.
The fourth term $\gT{\alpha\beta}$ is precomputed and stored in the verifier key, so no extra pairing is needed for it.
\end{tcolorbox}
 inside zero_knowledge.tex

%% ============================================================
\section{Groth16: A Pairing-Based SNARK}
%% ============================================================

Groth16 is a preprocessing SNARK for arithmetic circuit satisfiability. Its proof is only three group elements, verification is a single pairing equation, and it achieves perfect completeness and perfect zero knowledge. Soundness holds against adversaries that interact with the system via generic group operations.

We build the system in three layers: encode computation as a polynomial divisibility problem (Quadratic Arithmetic Program, QAP), construct an information-theoretic proof for it (Non-Interactive Linear Proof, NILP), then compile the NILP into group elements using pairings.

%% -----------------------------------------------------------
\subsection{Notation for Pairing Groups}
%% -----------------------------------------------------------

Let $(p, \G_1, \G_2, \GT, e, g, h)$ be an asymmetric bilinear group of prime order $p$. The generator of $\G_1$ is $g$, the generator of $\G_2$ is $h$, and $e : \G_1 \times \G_2 \to \GT$ is a bilinear pairing satisfying $e(g^a, h^b) = e(g,h)^{ab}$.

We represent group elements by their discrete logarithms using the bracket notation
\[
  \gone{a} \deq g^a \in \G_1, \qquad
  \gtwo{a} \deq h^a \in \G_2, \qquad
  \gT{a}   \deq e(g,h)^a \in \GT, \qquad a \in \Z_p.
\]
This notation makes linear algebra readable. For a known scalar $c$ or known matrix $M$:
\[
  \gone{a} + \gone{b} = \gone{a+b}, \qquad
  c \cdot \gone{a} = \gone{ca}, \qquad
  M \cdot \gone{\mathbf{a}} = \gone{M\mathbf{a}}.
\]
The pairing provides one level of multiplication across the two source groups:
\[
  \gone{a} \cdot \gtwo{b} \deq e\!\left(g^a, h^b\right) = \gT{ab}.
\]


%% -----------------------------------------------------------
\subsection{From Arithmetic Circuits to Quadratic Arithmetic Programs}
%% -----------------------------------------------------------

\subsubsection*{Arithmetic circuits and R1CS}

An arithmetic circuit over a finite field $\F_p$ consists of addition and multiplication gates over field elements. Any such circuit can be \emph{flattened}: every gate is rewritten as a single multiplicative constraint
\[
  \bigl(\textstyle\sum_i A_{q,i}\, a_i\bigr)
  \cdot
  \bigl(\textstyle\sum_i B_{q,i}\, a_i\bigr)
  =
  \textstyle\sum_i C_{q,i}\, a_i
  \qquad (q = 1,\ldots, n),
\]
where $z = (a_0, a_1, \ldots, a_m) \in \F_p^{m+1}$ is the \emph{assignment vector}: $a_0 = 1$ is a constant, $\phi = (a_1, \ldots, a_\ell)$ is the public statement, and $w = (a_{\ell+1}, \ldots, a_m)$ is the secret witness. The $n$ constraints are collected into three matrices $A, B, C \in \F_p^{n \times (m+1)}$ --- this is the \emph{rank-1 constraint system} (R1CS). Arithmetic circuit satisfiability is NP-complete: any NP statement can be compiled into an R1CS instance over a suitable field, so a proof system for R1CS is a proof system for all of NP.

\subsubsection*{Quadratic arithmetic programs}

Checking $n$ separate scalar equations one by one would be costly. The key algebraic trick is to compress them into a single polynomial divisibility test.

Fix $n$ distinct field elements $r_1, \ldots, r_n \in \F_p$, one for each constraint. Define the vanishing polynomial
\[
  t(X) = \prod_{q=1}^{n}(X - r_q).
\]
For each variable index $i \in \{0, \ldots, m\}$, interpolate three polynomials $u_i, v_i, w_i \in \F_p[X]$ of degree at most $n-1$ such that
\[
  u_i(r_q) = A_{q,i}, \qquad v_i(r_q) = B_{q,i}, \qquad w_i(r_q) = C_{q,i} \qquad \forall\, q.
\]
\todo{Add a visualisation showing how the columns of the $A$, $B$, $C$ matrices are read off and interpolated into the polynomials $u_i$, $v_i$, $w_i$.}
Given an assignment $z$, define the aggregated polynomials
\[
  U(X) = \sum_{i=0}^{m} a_i\, u_i(X), \qquad
  V(X) = \sum_{i=0}^{m} a_i\, v_i(X), \qquad
  W(X) = \sum_{i=0}^{m} a_i\, w_i(X).
\]
At each constraint point $r_q$, the R1CS equation $U(r_q)\cdot V(r_q) = W(r_q)$ must hold. Since $t(X)$ vanishes precisely at $r_1, \ldots, r_n$, this is equivalent to
\[
  \boxed{t(X) \;\Big|\; U(X)\cdot V(X) - W(X).}
\]
Equivalently, there exists a degree-$(n-2)$ quotient polynomial $h(X)$ such that
\[
  U(X)\cdot V(X) = W(X) + h(X)\,t(X).
\]
This is the \emph{quadratic arithmetic program} (QAP) relation for the circuit. The advantage over R1CS is that $n$ separate equations have been replaced by one polynomial identity checkable at a single random point --- exactly the Schwartz--Zippel trick.

%% -----------------------------------------------------------
\subsection{A Non-Interactive Linear Proof for QAP}
%% -----------------------------------------------------------

A \emph{non-interactive linear proof} (NILP) is an information-theoretic proof system where the prover's messages are linear functions of a common reference string $\sigma$. We now construct such a system for the QAP relation.

\todo{Explain why this NILP is non-interactive despite stemming from a Linear Interactive Proof (LIP): the CRS plays the role of the verifier's challenges, collapsing the interaction into a single message. Discuss the connection between 2-move LIPs and NILPs more carefully.}

\subsubsection*{Setup}

Pick a random trapdoor $\tau = (\alpha, \beta, \gamma, \delta, x) \in \F_p^{*5}$ and define the common reference string $\sigma$ as the collection of field evaluations
\[
\sigma =
\left(
\begin{aligned}
&\alpha,\; \beta,\; \gamma,\; \delta,\;\{x^i\}_{i=0}^{n-1},\\[2pt]
&\left\{\frac{\beta u_i(x)+\alpha v_i(x)+w_i(x)}{\gamma}\right\}_{i=0}^{\ell},\\[4pt]
&\left\{\frac{\beta u_i(x)+\alpha v_i(x)+w_i(x)}{\delta}\right\}_{i=\ell+1}^{m},\\[4pt]
&\left\{\frac{x^i\, t(x)}{\delta}\right\}_{i=0}^{n-2}
\end{aligned}
\right).
\]
The first block enables evaluation of arbitrary polynomials at $x$. The second block encodes the public witness terms, scaled by $\gamma$. The third encodes the private witness terms, scaled by $\delta$. The last encodes powers of $x$ weighted by $t(x)/\delta$, allowing the prover to compute $h(x)\,t(x)/\delta$ as a linear combination.

\todo{Discuss what $\alpha$, $\beta$, $\gamma$, and $\delta$ each do exactly: why they appear in the CRS, how they enforce consistency between $A$, $B$, $C$, and why the $\gamma$/$\delta$ separation is necessary.}

\subsubsection*{Prover}

Given a satisfying assignment $(a_1, \ldots, a_m)$, the prover computes $h(X)$ from the QAP relation, then picks two random blinding scalars $r, s \leftarrow \F_p$ and sets
\begin{align*}
  A &\deq \alpha + U(x) + r\delta, \\
  B &\deq \beta  + V(x) + s\delta, \\
  C &\deq \frac{\displaystyle\sum_{i=\ell+1}^{m} a_i\bigl(\beta u_i(x)+\alpha v_i(x)+w_i(x)\bigr) + h(x)\,t(x)}{\delta} + As + rB - rs\delta.
\end{align*}
The proof is $\pi = (A, B, C)$.

\subsubsection*{Verifier}

The verifier checks the single equation
\[
  A \cdot B
  \;\stackrel{?}{=}\;
  \underbrace{\alpha\beta}_{\text{(i)}}
  \;+\;
  \underbrace{\sum_{i=0}^{\ell} a_i \cdot \frac{\beta u_i(x)+\alpha v_i(x)+w_i(x)}{\gamma}}_{\text{(ii)}} \cdot\, \gamma
  \;+\;
  \underbrace{C \cdot \delta}_{\text{(iii)}}.
\]
The three terms on the right play distinct roles.
\begin{itemize}[noitemsep, leftmargin=2em]
  \item[(i)] $\alpha\beta$ is a fixed anchor from the setup. Its presence forces $A$ and $B$ to contain genuine $\alpha$ and $\beta$ offsets respectively, preventing a prover from constructing $A$ and $B$ independently of the structure of $\sigma$.
  \item[(ii)] The sum over $i = 0, \ldots, \ell$ incorporates the \emph{public inputs} $a_0, \ldots, a_\ell$ --- the statement $\phi$ that the verifier knows. These are precomputed in the CRS scaled by $\gamma$ and $\gamma$ cancels out, leaving the verifier to evaluate this term using only public data. Note that $\alpha$ and $\beta$ also appear here, but they are \emph{setup parameters} embedded in the precomputed CRS entries, not the circuit's inputs.
  \item[(iii)] $C \cdot \delta$ is where the private witness $a_{\ell+1}, \ldots, a_m$ is encoded. The $\delta$ scaling keeps it algebraically separated from term (ii), preventing the prover from mixing public and private contributions.
\end{itemize}

\subsubsection*{Simulator}

Given the trapdoor $\tau$, the simulator ignores the witness entirely. It picks $A, B \leftarrow \F_p$ uniformly at random and computes
\[
  C \deq \frac{A B \;-\; \alpha\beta \;-\; \displaystyle\sum_{i=0}^{\ell} a_i\bigl(\beta u_i(x)+\alpha v_i(x)+w_i(x)\bigr)}{\delta}.
\]

\subsubsection*{Why this achieves zero knowledge}

The blinding scalars $r$ and $s$ make $A$ and $B$ uniformly distributed in $\F_p$, independent of the witness. Given uniform $A$ and $B$, the verification equation uniquely determines $C$. The real distribution of $(A, B, C)$ is therefore identical to the simulated one: two uniform field elements followed by their forced companion. This is perfect zero knowledge.

\subsubsection*{Why this is sound against affine strategies}

An affine prover strategy computes $(A, B, C)$ as linear combinations of $\sigma$. Viewing the verification equation as a polynomial identity in the formal indeterminates $\alpha, \beta, \gamma, \delta, x$ and applying Schwartz--Zippel, any affine strategy that passes verification with non-negligible probability must satisfy the identity formally. A careful analysis of the monomials (particularly those involving $\alpha\beta$, $1/\gamma^2$, and $1/\delta^2$) forces the coefficients to encode a consistent assignment, from which the witness $(a_{\ell+1}, \ldots, a_m)$ can be extracted.

\todo{Add the full monomial analysis proof here, following Theorem~1 of the Groth (2016) paper.}

%% -----------------------------------------------------------
\subsection{Compiling to Pairings: The Generic Group Model}
%% -----------------------------------------------------------

The NILP works over a field, but we need a cryptographic argument --- the prover should not be able to read off $x$ from $\sigma$. The compilation step replaces every field element in $\sigma$ with its bracket encoding, and $(A, B, C)$ become $(\gone{A}, \gtwo{B}, \gone{C})$.

The compiled verification equation is:
\[
  \gone{A} \cdot \gtwo{B}
  \;=\;
  \gone{\alpha}\cdot\gtwo{\beta}
  \;+\;
  \sum_{i=0}^{\ell} a_i\,
    \left[\frac{\beta u_i(x)+\alpha v_i(x)+w_i(x)}{\gamma}\right]_{\!1}
    \cdot\, \gtwo{\gamma}
  \;+\;
  \gone{C}\cdot\gtwo{\delta}.
\]
The prover computes $\gone{A}$, $\gone{C}$, $\gtwo{B}$ as linear combinations of the CRS elements --- this is exactly the NILP prover running in the exponent. The verifier checks one pairing product equation using three pairings in total.

\textbf{Lifting soundness via the generic group model.} In the generic group model, a prover can only use the group elements in $\sigma$ through addition (and scalar multiplication) --- it cannot read their discrete logarithms. Any computation it performs on $\gone{\sigma_1}$ and $\gtwo{\sigma_2}$ is therefore a sequence of affine operations, corresponding exactly to choosing a proof matrix $\Pi$ and computing $\Pi\sigma$. This is precisely an affine strategy in the NILP sense. Combined with a \emph{disclosure-freeness} property of the CRS --- which guarantees that the structure of $\sigma$ cannot be exploited to choose a non-affine strategy --- the NILP's soundness against affine strategies lifts to soundness against all generic adversaries.

%% -----------------------------------------------------------
\subsection*{Summary: The Groth16 Proof System}
%% -----------------------------------------------------------

\begin{tcolorbox}[enhanced, colback=blue!3!white, colframe=blue!45!black,
                  arc=2.5pt, boxrule=1pt, left=6pt, right=6pt, top=5pt, bottom=5pt]
\small
\textbf{Setup.} Sample $\tau = (\alpha,\beta,\gamma,\delta,x) \leftarrow \F_p^{*5}$.
Publish $\sigma = \bigl([\sigma_1]_1,\, [\sigma_2]_2\bigr)$ as described above. Keep $\tau$ secret.

\medskip
\textbf{Prove.} Given $(a_1,\ldots,a_m)$ satisfying the circuit, pick $r,s \leftarrow \F_p$ and output
$\pi = \bigl(\gone{A},\, \gone{C},\, \gtwo{B}\bigr) \in \G_1^2 \times \G_2$.

\medskip
\textbf{Verify.} Check:
$\gone{A} \cdot \gtwo{B} = \gone{\alpha}\cdot\gtwo{\beta} + \displaystyle\sum_{i=0}^{\ell} a_i\,\gone{(\cdot)_i}\cdot\gtwo{\gamma} + \gone{C}\cdot\gtwo{\delta}$.

\medskip
\textbf{Proof size:} \quad $2$ elements in $\G_1$, \quad $1$ element in $\G_2$ \quad $= 3$ group elements.\\
\textbf{Verification:} \quad $1$ pairing product equation requiring $3$ online pairings:
$e(\gone{A}, \gtwo{B})$,\; $e\!\left(\textstyle\sum a_i K_i,\, \gtwo{\gamma}\right)$,\; $e(\gone{C}, \gtwo{\delta})$.
The fourth term $\gT{\alpha\beta}$ is precomputed and stored in the verifier key, so no extra pairing is needed for it.
\end{tcolorbox}


\end{document}
